% =====================================================================
%  Chapter 27 review -- the 14-station revision contract.
%  Revision-only: tables, one-liners, pointers back via \secref/\thmref/
%  \defref/\figref/\coderef/\probref.  No re-derivations.
%  Summit: CF 1585E "Frequency Queries" (2400, L6) -- a rooted tree with an
%  integer a_i on each node and q offline questions (v,l,k): read the values on
%  the path v -> root, tally each value's frequency, keep values with freq >= l,
%  order the survivors by frequency INCREASING (ties arbitrary), and report the
%  k-th value (or -1).  This is the chapter's own "non-invertible offline
%  root-path statistic": a plain DFS from the root is the window (push a value on
%  entry, pop it on exit -- the exact O(1) add/remove primitive of 27.1, now
%  riding the recursion stack instead of Mo's two pointers), cnt[val] + a
%  frequency-of-frequency Fenwick select the k-th, and a per-frequency value
%  stack returns the answer value.  O((n+q) log n).  Reuses TWO chapter tools
%  (the Mo frequency table + the offline batch answering of Tarjan) and ONE
%  earlier-chapter tool (the Ch.5 depth/parent DFS).  Both official samples
%  hand-replayed BEFORE writing: [1 -1 1] and [1 1 2 1 -1].
%  Exports prob:c21-summit, lst:c21-summit, lst:c21-kit (kit lives in 27.5).
%  Mixed set EXACTLY 10, all live-verified this session (names/ratings/tags
%  re-pulled from the CF problemset API), un-owned, ASCENDING by rating; each
%  hint routes via the signature (def:c21-signature) or REJECTS Mo and names the
%  true tool (Fenwick/Euler, HLD, DSU-on-tree, centroid, LCA).  The canonical
%  tree-Mo problems 375D (Ch.17), 258E, 86D are owned elsewhere and NEVER
%  re-linked here.
% =====================================================================

\chaptersection{Chapter 27 Summary}

One sentence: \emph{when one fixed tree faces a huge offline batch of subtree or
path queries whose statistic has \textbf{no inverse} -- distinct colours,
\(\sum\mathrm{freq}^2\), "how many values appear \(\ge l\) times" -- you stop
being online and instead \emph{reorder} the queries so consecutive ones differ by
a single node, maintaining the statistic with an \(O(1)\) add/remove across a
\(\bigO{(n+q)\sqrt n}\) block-order sweep; and when the queries only want
\(\lca\), Tarjan's one-DFS \(+\) DSU answers the whole batch in
\(\bigO{(n+q)\alpha}\) (\(\secref{c21-moengine}\), \(\thmref{c21-complexity}\),
\(\thmref{c21-tarjan}\)).} Chapter~26 handed you the "many overlapping offline
node-set queries, non-invertible statistic" case at its method boundary; this
chapter is the answer to it. Each row below is a pointer, not a proof.

\begin{center}
\renewcommand{\arraystretch}{1.25}
\setlength{\tabcolsep}{3pt}
\begin{tabular}{@{}p{0.22\linewidth} p{0.53\linewidth} p{0.19\linewidth}@{}}
\toprule
\textbf{Result} & \textbf{Statement} & \textbf{Where} \\
\midrule
Block order & Sort \(q\) offline ranges by \((\lfloor L/B\rfloor,R)\), creep one
window \([curL,curR]\) with an \(O(1)\) \texttt{add}/\texttt{rm}; with
\(B=\Theta(n/\sqrt q)\) the total pointer movement is \(\bigO{(n+q)\sqrt n}\). &
\defref{c21-mo}, \thmref{c21-complexity} \\
XOR engine & Prefix-xor turns "subarray xor \(=k\)" into "count pairs with
\(px[p]\oplus px[q]=k\) in \([l-1,r]\)"; insert \(v\) forms \(cnt[v\oplus k]\) new
pairs -- Mo on \(px[0..n]\). & \probref{c21-xor}, \coderef{c21-xor} \\
Subtree flatten & subtree\((v)\) is the contiguous preorder range
\([\tin(v),\tout(v)]\); run the same sweep on it for a non-invertible subtree
statistic. & \defref{c21-flatten}, \thmref{c21-subtree-range} \\
Root-path stack & CF 463E's queried set is the current root path: a plain DFS
 carries the window -- per-prime \(\mathit{last}[p]\) with \(O(1)\) push/undo,
 one rebuild per block of the \(\le50\) updates. \emph{No Mo, no lifting.} & \probref{c21-caisa}, \coderef{c21-caisa} \\
Path toggle & A doubled tour writes each node at entry and exit; crossing an
occurrence \textbf{toggles} membership, and a path is
\([st(u),st(v)]\) if \(u\) ancestors \(v\), else \([en(u),st(v)]\) \emph{plus the
LCA at read time}. & \defref{c21-toggle}, \thmref{c21-pathsplit} \\
Tarjan LCA & One DFS, three colours; union a returned child, stamp the parent; a
query fires when its second endpoint turns BLACK, reading
\(\mathit{anc}[\mathit{find}(\text{partner})]\). \(\bigO{(n+q)\alpha}\), offline only. &
\defref{c21-tarjan}, \thmref{c21-tarjan} \\
The pick & Offline + overlapping + non-invertible + \(O(1)\) update \(\Rightarrow\)
Mo; else Fenwick (invertible), HLD/LCT (online), DSU-on-tree (subtree-all),
virtual tree (per-set), centroid (all-pairs). & \defref{c21-signature}, \defref{c21-pick} \\
\bottomrule
\end{tabular}
\end{center}

\medskip
{\small\textbf{Algorithms from this chapter.}}
\begin{center}
\renewcommand{\arraystretch}{1.2}
\setlength{\tabcolsep}{3pt}
\begin{tabular}{@{}p{0.36\linewidth} l l >{\raggedright\arraybackslash}p{0.20\linewidth}@{}}
\toprule
\textbf{Algorithm} & \textbf{Prep} & \textbf{Per run} & \textbf{Code} \\
\midrule
Generic Mo block-order sweep & \(\bigO{q\log q}\) sort & \(\bigO{(n+q)\sqrt n}\) & \coderef{c21-mo} \\
CF 617E: Mo over prefix xors & \(\bigO{n}\) prefix & \(\bigO{(n+m)\sqrt n}\) & \coderef{c21-xor} \\
Subtree flatten + Mo on ranges & \(\bigO{n}\) DFS & \(\bigO{(n+q)\sqrt n}\) & \coderef{c21-flatten} \\
CF 463E: DFS stack + last-seen primes & \(\bigO{A}\) sieve & \(\bigO{51\,n\cdot 7+q}\) & \coderef{c21-caisa} \\
Doubled-tour path toggle + LCA split & \(\bigO{n}\) tour & \(\bigO{(n+q)\sqrt n}\) & \coderef{c21-euler-toggle} \\
Tarjan offline LCA & \(\bigO{n}\) DFS+DSU & \(\bigO{q\alpha}\) & \coderef{c21-tarjan} \\
Recognition gate / boundary pick & -- & -- & \defref{c21-signature}, \defref{c21-pick} \\
Mo kit (one decision listing) & -- & -- & \coderef{c21-kit} \\
CF 1585E (Summit): root-path freq-of-freq & \(\bigO{n}\) build & \(\bigO{(n+q)\log n}\) & \coderef{c21-summit} \\
\bottomrule
\end{tabular}
\end{center}

\chaptersection{Master Decision Flowchart}

Post-study ladder: walk top to bottom and stop at the first branch that matches.
It decides \emph{whether the problem is a Mo-on-trees job at all} (offline batch,
overlapping subtree/path sets, non-invertible statistic, \(O(1)\) add/remove) and,
if so, \emph{which array} the sweep runs on (preorder flatten for subtrees, the
doubled tour for paths). Everything that fails the top gate is handed to a
cheaper or online engine: invertible ranges to Euler \(+\) Fenwick (\(\chref{c24}\)),
online paths to HLD (\(\chref{c16}\)) / link-cut (\(\chref{c23}\)), per-set folds
to virtual trees (\(\chref{c30}\)), whole-tree distance shapes to centroid
(\(\chref{c18}\)), and pure LCA batches to Tarjan (\(\secref{c21-tarjan}\)).

\begin{lstlisting}[language={}, basicstyle=\ttfamily\small, numbers=none, frame=single]
  Are ALL queries known up front, over ONE fixed tree, with NO updates between them?
  |
  |-- An update or answer must come BETWEEN queries (online / interleaved)?
  |     +-- YES: NOT Mo -- block order needs the whole batch.
  |         Online path range -> HLD (Ch.22, c16); dynamic tree -> LCT (Ch.25, c23);
  |         online subtree range -> Euler tour + segment tree (Ch.5, c24).
  |
  |-- Does the statistic telescope: answer = PREF(r) - PREF(l-1)?
  |     +-- YES: NOT Mo -- pay a LOG, not a SQRT.
  |         Sum / xor / depth-parity-normalised add -> Fenwick over the flatten
  |         (Ch.5, c24).  Template: Ch.18's CF 383C (parity BIT, c14).
  |
  |-- Queries ask ONLY for lca(u,v), offline, and you want it near-linear?
  |     +-- YES: Tarjan DFS+DSU, O((n+q) alpha), no lifting table (27.4, lst:c21-tarjan).
  |         (If you ALSO need online LCAs, keep Ch.16 lifting.)
  |
  |-- Each query a SMALL marked set to fold independently, sum k << n?
  |     +-- YES: virtual tree (Ch.26, c30) -- the sets do not overlap as ranges.
  |
  |-- No query list: a whole-tree pair / distance distribution (sum over all paths)?
  |     +-- YES: centroid decomposition (Ch.24, c18).
  |
  |-- Subtree-only, statistic merges small-to-large, query set = (most) nodes?
  |     +-- YES: price DSU on tree (Ch.17, c28) -- O(n log n) can beat O((n+q) sqrt n).
  |
  |-- STILL HERE: offline, overlapping subtree/path ranges, non-invertible O(1) statistic
  |     -> THIS chapter.  Choose the array the sweep runs on:
  |        subtrees -> preorder flatten P, Mo on [tin(v), tout(v)]       (27.2, lst:c21-flatten)
  |        paths    -> doubled tour E, TOGGLE, split at lca + add LCA    (27.3, lst:c21-euler-toggle)
  |        arrays   -> block order (L/B, R), O(1) add/rm                 (27.1, lst:c21-mo)
\end{lstlisting}

\chaptersection{Common Mistakes Checklist}

\begin{itemize}[nosep, leftmargin=1.4em]
  \item \textbf{Paying the \(\sqrt n\) tax on an invertible statistic.} If the
        answer is \(\mathrm{pref}(r)-\mathrm{pref}(l-1)\), a Fenwick beats Mo by a
        \(\sqrt n/\log n\) factor.  Ch.~18's CF 383C \emph{looks} like Mo but telescopes
        after the \((-1)^{\depth}\) normalisation (\(\chref{c14}\)).
  \item \textbf{Forgetting the empty prefix \(px[0]\) in CF 617E.} A subarray from
        the left edge pairs \(px[l-1]\) with a later prefix; if \(L=l-1\) never
        adds index \(0\), every left-edge subarray is under-counted
        (\(\coderef{c21-xor}\)).
  \item \textbf{\texttt{rm} updating in the wrong order.} \texttt{rm} must
        \emph{decrement} the frequency \emph{before} subtracting the formed pairs,
        or it removes a pair that is itself leaving (\(\coderef{c21-xor}\)).
  \item \textbf{Sorting queries by \((R,\text{block of }L)\) instead of the reverse.}
        The bound needs one pointer monotone per group and the other confined to
        width \(B\); swapping the keys gives \(O(nq)\) again
        (\(\thmref{c21-complexity}\), \(\defref{c21-mo}\)).
  \item \textbf{Missing the LCA in path Mo (case ii).} When \(\lca(u,v)\neq u\)
        the window represents the path \emph{minus} the LCA; forget the read-time
        \texttt{insert(lca)} and answers go wrong exactly when the LCA carries a
        new colour (\(\thmref{c21-pathsplit}\), walk 3 of \(\secref{c21-path}\)).
  \item \textbf{Double-adding the LCA in case (i).} When \(u\) is the ancestor the
        LCA is \emph{already} oddly represented; toggling it again removes it.
        The extra insert is for case (ii) only (\(\coderef{c21-euler-toggle}\)).
  \item \textbf{Attaching a Tarjan query to only one endpoint.} It fires at the
        \emph{second} endpoint's death; if the first endpoint is where it lives, it
        checks a WHITE/GREY partner and is silently never answered
        (\(\defref{c21-tarjan}\), \(\coderef{c21-tarjan}\)).
  \item \textbf{Running Mo with interleaved updates.} Block order is defined only
        for a fixed array; a point update between queries invalidates the whole
        sweep -- that is the HLD / sqrt-decomposition-over-time contract
        (\(\defref{c21-signature}\), \(\secref{c21-boundary}\)).
\end{itemize}

\chaptersection{WA Clinic}

Verdicts are symptoms; the causes are the short list this chapter drilled. Diagnose
top-down -- every leaf names a listing.

\begin{figure}[H]
\centering
\begin{tikzpicture}[diagram,
    v/.style={boxnode, fill=primaryblue!12, draw=primaryblue, text width=22mm, align=center},
    b/.style={boxnode, fill=difforange!12, draw=difforange, text width=12mm},
    l/.style={boxnode, fill=successgreen!10, draw=successgreen, text width=48mm,
              font=\scriptsize, align=left}]
  \node[v] (verdict) {\textbf{Verdict on the judge}};
  \node[b, below left=8mm and 5mm of verdict] (wa) {\textbf{WA}};
  \node[b, below right=8mm and 5mm of verdict] (tle) {\textbf{TLE / RE}};
  \node[l, below right=1mm and 2mm of wa.south] (wa1)
    {path-Mo case (ii) read without the extra LCA, or case (i) toggled the LCA twice --
     wrong exactly on the queries where the LCA adds a colour (\(\thmref{c21-pathsplit}\), \(\coderef{c21-euler-toggle}\))};
  \node[l, below left=1mm and 2mm of wa.south] (wa2)
    {617E dropped \(px[0]\) or ran \texttt{rm} before decrementing the count -- off-by-pairs
     on left-edge subarrays (\(\coderef{c21-xor}\)); or 463E answered after pushing its
     own primes -- every match degenerates to \(v\) itself (\(\coderef{c21-caisa}\))};
  \node[l, below right=1mm and 2mm of tle.south] (tle1)
    {Mo on an invertible statistic (a Fenwick would be \(\bigO{\log n}\)) or block size
     \(B\) set to a constant far from \(n/\sqrt q\) -- the sweep degrades (\(\thmref{c21-complexity}\), \(\defref{c21-pick}\))};
  \node[l, below left=1mm and 2mm of tle.south] (tle2)
    {Mo fired on an online / update-interleaved query (block order illegal), or a recursion
     DFS overflowed at \(n=2\cdot10^5\) -- go HLD / iterative (\(\defref{c21-signature}\), \(\secref{c21-tarjan}\))};
  \draw[flow] (verdict) -- (wa);
  \draw[flow] (verdict) -- (tle);
  \draw[flow] (wa) -- (wa1);
  \draw[flow] (wa) -- (wa2);
  \draw[flow] (tle) -- (tle1);
  \draw[flow] (tle) -- (tle2);
\end{tikzpicture}
\caption*{WA: the path Mo LCA handled wrong, or the prefix / parity normalisation
mis-set. TLE or RE: paying \(\sqrt n\) where \(\log n\) applies, a badly chosen block
size, or Mo on an online problem (and a deep recursion at \(n=2\cdot10^5\)).}
\end{figure}

\chaptersection{Quick Reference Card}

\begin{quickref}[Chapter 27 at a glance]
\textbf{Mo on trees \(=\) reorder offline subtree/path queries so neighbours differ
by one node, then sweep a window with an \(O(1)\) add/remove; all \(q\) answers in
\(\bigO{(n+q)\sqrt n}\). For \(q\) LCAs only, Tarjan's DFS \(+\) DSU does the batch
in \(\bigO{(n+q)\alpha}\).}

\medskip
\textbf{Engine (27.1).} Block size \(B=n/\sqrt q\); sort by
\((\lfloor L/B\rfloor,R)\); creep \([curL,curR]\) with \texttt{add}/\texttt{rm};
read the statistic. CF 617E: window over prefix-xors, insert \(v\) adds
\(cnt[v\oplus k]\) pairs (\(\thmref{c21-complexity}\), \(\coderef{c21-mo}\),
\(\coderef{c21-xor}\)).

\medskip
\textbf{Subtrees (27.2).} Flatten: \(P[\tin(v)]=v\), subtree\((v)=[\tin(v),\tout(v)]\);
Mo on those ranges for a distinct/freq statistic. Invertible signed updates skip Mo
(Ch.~18's CF 383C parity BIT, \(\chref{c14}\)); root-path batches with \(\le50\)
updates ride the DFS stack: per-prime \(\mathit{last}[p]\) push/undo, one rebuild per
block (\(\defref{c21-flatten}\), \(\thmref{c21-subtree-range}\),
\(\coderef{c21-caisa}\)).

\medskip
\textbf{Paths (27.3).} Doubled tour \(E[1..2n]\); toggling gives represented\(=\)odd
occurrences; window \([st(u),st(v)]\) if \(u\) ancestors \(v\), else \([en(u),st(v)]\)
\emph{plus the LCA at read time} (\(\defref{c21-toggle}\), \(\thmref{c21-pathsplit}\),
\(\coderef{c21-euler-toggle}\)).

\medskip
\textbf{Tarjan (27.4).} WHITE/GREY/BLACK DFS; union a returned child,
\(\mathit{anc}[\mathit{find}(u)]=u\); when \(u\) turns BLACK answer its queries whose
partner is BLACK via \(\mathit{anc}[\mathit{find}(\text{partner})]\)
(\(\thmref{c21-tarjan}\), \(\coderef{c21-tarjan}\)).

\medskip
\textbf{The pick (27.5).} Offline \emph{and} overlapping \emph{and} non-invertible
\emph{and} \(O(1)\) update \(\Rightarrow\) Mo; invertible \(\to\) Fenwick
(\(\chref{c24}\)); online \(\to\) HLD/LCT (\(\chref{c16},\chref{c23}\)); per-set
\(\to\) virtual tree (\(\chref{c30}\)); subtree-all \(\to\) DSU-on-tree
(\(\chref{c28}\)); all-pairs \(\to\) centroid (\(\chref{c18}\)); LCA-batch \(\to\)
Tarjan (\(\defref{c21-signature}\), \(\defref{c21-pick}\)).
\end{quickref}

\chaptersection{Write-From-Memory Drills}

Five signature skeletons, each lifted from a listing already in the chapter. Cover the
marked line, reproduce from memory, then check the answer one line below.

\begin{lstlisting}[caption={Drill 1 -- the Mo block-order sweep (from lst:c21-mo)}]
sort(queries, cmp);                    // cmp: block of L, then R
curL = 0; curR = -1;                   // start EMPTY
for (qu : queries){
    while (curL > qu.l) add(--curL);   while (curR < qu.r) add(++curR);
    while (curL < qu.l) rm(curL++);    while (curR > qu.r) rm(curR--);
    ans[qu.idx] = cur;
}
\end{lstlisting}
\noindent{\scriptsize Key: grow first, then shrink, in that order keeps the window
valid; \texttt{add}/\texttt{rm} must be \(O(1)\) and mutually inverse; total movement
\(\bigO{(n+q)\sqrt n}\) with \(B=n/\sqrt q\) (\(\thmref{c21-complexity}\)).}

\begin{lstlisting}[caption={Drill 2 -- CF 617E add/rm over prefix xors (from lst:c21-xor)}]
void add(int t){ int v=px[t]; cur += cnt[v ^ k]; cnt[v]++; }   // pairs formed, then store
void rm (int t){ int v=px[t]; cnt[v]--;  cur -= cnt[v ^ k]; }   // un-store, then pairs lost
// query (l,r) -> window [l-1, r] over px[0..n]; remember px[0]=0
\end{lstlisting}
\noindent{\scriptsize Key: a subarray xor is \(px[q]\oplus px[p]\); inserting \(v\)
creates \(cnt[v\oplus k]\) good pairs; \texttt{rm} decrements \emph{first} so the
leaving element is not counted against itself; \(cur\) is \code{long long}
(\(\probref{c21-xor}\)).}

\begin{lstlisting}[caption={Drill 3 -- subtree flatten + Mo range (from lst:c21-flatten)}]
void dfs(int u,int p){ tin[u]=++t; P[t]=u;
    for(int v:adj[u]) if(v!=p){ dep[v]=dep[u]+1; dfs(v,u);} tout[u]=t; }
// subtree(v) aggregate == Mo over array range [ tin[v] , tout[v] ] on P
//   add(pos) inserts node P[pos] into the statistic;  rm(pos) erases it.
\end{lstlisting}
\noindent{\scriptsize Key: preorder makes every subtree one contiguous block, so the
array engine runs unchanged on \([\tin(v),\tout(v)]\); only \texttt{add} now touches
\(P[\text{pos}]\) (\(\defref{c21-flatten}\), \(\thmref{c21-subtree-range}\)).}

\begin{lstlisting}[caption={Drill 4 -- path window + toggle (from lst:c21-euler-toggle)}]
if (st[u] > st[v]) swap(u,v);                  // keep st[u] < st[v]
if (is_ancestor(u,v)){ L=st[u]; R=st[v]; extraLca=false; }
else                 { L=en[u]; R=st[v]; extraLca=true; }   // window misses the lca
// add(pos): x=eu[pos]; held[x]^=1; (held[x]? insert(x): erase(x));   rm = add
// if extraLca:  insert(lca(u,v)); read cur; erase(lca(u,v));
\end{lstlisting}
\noindent{\scriptsize Key: a path is a doubled-tour window whose represented set is the
odd-occurrence nodes; when \(u\) is the ancestor the LCA is already inside, otherwise it
is \emph{zero} occurrences and must be added at read time (\(\thmref{c21-pathsplit}\)).}

\begin{lstlisting}[caption={Drill 5 -- Tarjan offline LCA core (from lst:c21-tarjan)}]
void dfs(int u,int p){
  col[u]=1; fa[u]=u; anc[u]=u;
  for(v : adj[u]) if(v!=p){ dfs(v,u); fa[find(v)]=find(u); anc[find(u)]=u; }
  for(Q : qlist[u]) if(col[Q.other]==BLACK) ans[Q.id]=anc[find(Q.other)];
  col[u]=BLACK;
}   // attach every query (u,v) to BOTH u and v
\end{lstlisting}
\noindent{\scriptsize Key: after each child returns its finished region merges up and
re-stamps \(u\); the query fires at the second endpoint's death, reading the stamp of
the partner's set -- the merge point of the two paths (\(\thmref{c21-tarjan}\)).}

\chaptersection{Complete Problem Checklist}

Every problem this chapter offers -- two fully solved in-section, worked traces on the
running example \(R\), and one Summit on its own CF samples. Tick \emph{only} after an
unaided solve; a hinted attempt counts as half.

\begin{center}
\footnotesize
\setlength{\tabcolsep}{4pt}
\renewcommand{\arraystretch}{1.25}
\begin{tabular}{@{}c p{0.16\linewidth} l p{0.20\linewidth} >{\raggedright\arraybackslash}p{0.26\linewidth} c@{}}
\toprule
\textbf{\#} & \textbf{Problem} & \textbf{Source} & \textbf{Key technique} & \textbf{What to look for} & \textbf{Done?} \\
\midrule
\multicolumn{6}{l}{\textbf{Section 27.1 -- the Mo engine (array)}} \\
W27.1 & block order on \(n{=}9,B{=}3\) & worked & \lv{3} sort \((L/B,R)\), creep & 13 moves; \(R\) monotone per block, \(L\) jumps \(\le B\) & \(\square\) \\
W27.1b & XOR and Favorite Number & CF 617E & \lv{4} Mo over prefix xors & \(px[0]\) counted; \texttt{rm} decrements first; samples \(\to\) 7,0 and 9,4,4 & \(\square\) \\
\midrule
\multicolumn{6}{l}{\textbf{Section 27.2 -- subtrees}} \\
W27.2 & subtree\((2)\) of \(R\) as a range & worked & \lv{3} preorder flatten & \(\subtree(2)=P[2..5]=\{2,4,8,5\}\); three distinct colours & \(\square\) \\
W27.2b & Caisa and Tree & CF 463E & \lv{4} DFS stack + last-seen primes & read before push (strict ancestor); \(\le50\) updates \(\to\) 51 rebuilds; sample \(\to\) -1,1,2,-1,1 & \(\square\) \\
\midrule
\multicolumn{6}{l}{\textbf{Section 27.3 -- paths (worked only)}} \\
W27.3 & paths \((4,8),(8,6),(2,3)\) on \(R\) & worked & \lv{5} toggle + LCA split & \((4,8)\to2\); \((8,6)\to3\) with LCA; \((2,3)\to3\) vs lying 2 & \(\square\) \\
\midrule
\multicolumn{6}{l}{\textbf{Section 27.4 -- Tarjan (worked only)}} \\
W27.4 & five LCAs on \(R\) & worked & \lv{4} DFS + DSU + colours & union-on-return; fire at 2nd BLACK; answers \(1,2,2,3,1\) & \(\square\) \\
\midrule
\multicolumn{6}{l}{\textbf{Section 27.5 -- recognition / boundary (worked only)}} \\
W27.5 & three specs, three verdicts & worked & \lv{5} the four tests & path-distinct \(\to\) Mo; subtree-all \(\to\) DSU; updates \(\to\) reject & \(\square\) \\
\midrule
\multicolumn{6}{l}{\textbf{Summit (chapter end)}} \\
W27 & Frequency Queries & CF 1585E & \lv{6} root-path freq-of-freq & DFS stack window + BIT k-th; samples \(\to\) 1 -1 1 / 1 1 2 1 -1 & \(\square\) \\
\bottomrule
\end{tabular}
\end{center}

\chaptersection{Recommended Study Order}

Four focused days. Day~1 the engine and CF 617E; day~2 the flatten, CF 463E and its
contrast to Mo; day~3 the doubled-tour toggle and Tarjan; day~4 the boundary pick, the
Summit and the mixed set.

\begin{center}
\small
\renewcommand{\arraystretch}{1.25}
\setlength{\tabcolsep}{3pt}
\begin{tabular}{@{}llp{0.55\linewidth}@{}}
\toprule
\textbf{Day} & \textbf{Focus} & \textbf{Tasks} \\
\midrule
1 & Read + do 27.1 & run the block order on \(n{=}9,B{=}3\) to 13 moves (\(\coderef{c21-mo}\)); prove \(\bigO{(n+q)\sqrt n}\) (\(\thmref{c21-complexity}\)); then CF 617E, replaying the prefix-xor window to \(7,0\) and \(9,4,4\) (\(\probref{c21-xor}\)) \\
2 & Read + do 27.2 & flatten subtree\((2)\) of \(R\) and read its three colours (\(\coderef{c21-flatten}\)); then CF 463E, carrying \(\mathit{last}[p]\) on the DFS stack and replaying to \(-1,1,2,-1,1\) (\(\probref{c21-caisa}\)) -- then say out loud why it needed \emph{neither} Mo nor lifting \\
3 & Read + do 27.3, 27.4 & walk the three path windows and the LCA-parity trap (\(\thmref{c21-pathsplit}\)); run the five-query Tarjan log to \(1,2,2,3,1\) (\(\thmref{c21-tarjan}\)) \\
4 & Read + do 27.5, Summit & place the three specs on the pick (\(\defref{c21-pick}\)); then CF 1585E, replaying both samples (\(\probref{c21-summit}\)); finish with the first 3 mixed entries \\
\bottomrule
\end{tabular}
\end{center}

\chaptersection{Difficulty Progression}

Tiers are competency gates -- each row names the \emph{skill} the tier unlocks.

\begin{center}
\small
\renewcommand{\arraystretch}{1.3}
\setlength{\tabcolsep}{3pt}
\begin{tabular}{@{}lp{0.30\linewidth}p{0.42\linewidth}@{}}
\toprule
\textbf{Tier} & \textbf{Problems} & \textbf{What passing this tier proves} \\
\midrule
\lv{3} \ \(\bigstar\bigstar\bigstar\)
  & block order by hand; subtree flatten
  & you can creep a window, count pointer moves, and read \([\tin(v),\tout(v)]\) \\
\lv{4} \ \(\bigstar\bigstar\bigstar\bigstar\)
  & CF 617E, CF 463E, Tarjan batch
  & you code an \(O(1)\) \texttt{add}/\texttt{rm}, carry a push/undo table on the DFS stack, and fire a Tarjan query at the right death \\
\lv{5} \ \(\bigstar\bigstar\bigstar\bigstar\bigstar\)
  & path toggle + LCA split, the boundary pick
  & you represent a path as a doubled-tour window, handle the parity trap, and \emph{reject} Mo to the right tool \\
\lv{6} \ \(\bigstar\bigstar\bigstar\bigstar\bigstar\bigstar\)
  & CF 1585E (Summit), mixed at CF~2300--2400
  & you recast the DFS stack as the Mo window and select a k-th-by-frequency answer with a frequency-of-frequency structure \\
\bottomrule
\end{tabular}
\end{center}

\chaptersection{Reduction Gallery}

Every entry is ``one question shape \(\to\) one move (or an honest hand-off).'' Each
line collapses a problem family onto a listing already in the chapter.

\begin{itemize}[nosep, leftmargin=1.4em]
  \item ``count subarrays with \(\mathit{xor}=k\), offline'' \(\to\) prefix-xor pairs,
        Mo window over \(px[0..n]\), \(O(1)\) insert/delete
        (\(\probref{c21-xor}\), \(\coderef{c21-xor}\)).
  \item ``distinct colours / \(\sum\mathrm{freq}^2\) over subtree\((v)\), offline''
        \(\to\) preorder flatten + Mo on \([\tin(v),\tout(v)]\)
        (\(\thmref{c21-subtree-range}\), \(\coderef{c21-flatten}\); CF 375D is the
        sibling -- its URL lives in Ch.17, \chref{c28}).
  \item ``add with sign flipping by depth, read a node'' (CF 383C) \(\to\) multiply
        every stored quantity by \((-1)^{\depth}\), one difference BIT
        (solved in Ch.~18, \(\chref{c14}\)).
  \item ``deepest ancestor sharing a factor, rare updates'' (CF 463E) \(\to\) the
        DFS stack \emph{is} the root-path window; per-prime last-seen push/undo,
        one rebuild per static block (\(\probref{c21-caisa}\), \(\coderef{c21-caisa}\)).
  \item ``distinct colours on path \(u\)-\(v\), offline'' \(\to\) doubled-tour toggle
        window + read-time LCA insert (\(\thmref{c21-pathsplit}\),
        \(\coderef{c21-euler-toggle}\)).
  \item ``answer \(q\) LCAs offline, no lifting table'' \(\to\) Tarjan DFS \(+\) DSU,
        fire at the second endpoint's death (\(\thmref{c21-tarjan}\),
        \(\coderef{c21-tarjan}\)).
  \item ``frequency histogram of the root-to-\(v\) path, k-th by count'' (CF 1585E)
        \(\to\) the DFS stack \emph{is} the window; cnt[val] + freq-of-frequency BIT
        (\(\probref{c21-summit}\), \(\coderef{c21-summit}\)).
  \item ``invertible range sum/xor on a subtree or path'' \(\to\) \emph{not} Mo:
        Euler \(+\) Fenwick / segment (\(\chref{c24}\), \(\defref{c21-pick}\)).
  \item ``online path range with updates'' \(\to\) \emph{not} Mo: HLD (\chref{c16})
        or link-cut (\chref{c23}); ``per-set fold'' \(\to\) virtual tree (\chref{c30}).
  \item ``subtree statistic for every node, mergeable'' \(\to\) price DSU-on-tree
        (\chref{c28}); ``all-pairs distance'' \(\to\) centroid (\chref{c18})
        (\(\secref{c21-boundary}\)).
\end{itemize}

\chaptersection{Mixed Problem Set}

Ten entries, each live-verified against the current Codeforces problem set (names,
ratings and tags re-pulled from the \texttt{problemset.problems} API this session) and
absent from every other roster in the book (the canonical tree-Mo problems CF 375D in
Ch.~17, CF 258E and CF 86D are owned elsewhere and not re-linked here). None says
``use Mo''; several deliberately \emph{do not} want it, and the drill is to run them
through the four tests (\(\defref{c21-signature}\)) and say so. Each hint names the
reduction (or the rejection) and the listing it leans on, ordered ascending by rating.

\begin{mixedset}
  \practice{Alyona and a tree}{CF 739B}
    {https://codeforces.com/problemset/problem/739/B}{CF 1900 L4}
    {Tags \emph{binary search, data structures, dfs, graphs, trees}. For each node \(v\),
     count ancestors \(u\) with \(\dist(u,v)\le a_v\). This is a \emph{root-path} sweep --
     the DFS stack \emph{is} the window (the Summit's shape, \(\probref{c21-summit}\)) --
     but the count telescopes along the depth axis: binary-lift to the highest valid
     ancestor, then a difference/BIT over depth. So it lives in world~1 (invertible,
     \(\secref{c21-subtree}\)) rather than Mo: \(\bigO{n\log n}\), not
     \(\bigO{(n+q)\sqrt n}\). Price both; the honest answer is the depth-difference.}
  \practice{Tree Queries}{CF 1328E}
    {https://codeforces.com/problemset/problem/1328/E}{CF 1900 L4}
    {Each query names \(a,b,x,y\): does \emph{some} path from \(a\) to \(b\) use the tree
     edge \((x,y)\)? Root once, orient the edge to its deeper end \(c\); the edge is on
     the \(a\)-\(b\) path iff exactly one of \(a,b\) lies inside subtree\((c)\) -- two
     \(\tin/\tout\) containment tests. Online per query and prefix-shaped: test~1 fails
     and test~3 fails (\(\defref{c21-signature}\)); reject Mo, keep
     \(\chref{c24}+\chref{c11}\). The same containment arithmetic 27.2's flatten hands
     over for free.}
  \practice{Passable Paths (easy version)}{CF 1702G1}
    {https://codeforces.com/problemset/problem/1702/G1}{CF 1900 L4}
    {Each query names a \emph{set} of vertices: is one simple path passable through all
     of them? Pick the set's two farthest vertices (diameter by \(\dist\) via
     \(\lca\)) and test every \(w\) for \(\dist(x,w)+\dist(w,y)=\dist(x,y)\) -- a
     membership test in pure distance arithmetic (\(\thmref{c11-distf}\)). Queries are
     independent checks, not overlapping ranges to amortise (test~2 empty) and nothing
     histogram-like is carried (test~3 fails). Reject; go to LCA (\(\chref{c11}\)).}
  \practice{Number of Simple Paths}{CF 1454E}
    {https://codeforces.com/problemset/problem/1454/E}{CF 2000 L4}
    {A connected graph on \(n\) vertices and \(n\) edges -- exactly one cycle -- count
     ordered pairs joined by \emph{exactly one} simple path. There is \emph{no query
     list} at all (test~2 empty): the answer is one global count from the cycle's
     attached subtree sizes, a parity-free inclusion-exclusion over the cycle.
     Structurally the cleanest rejection in the set -- count by structure, never sweep
     (\(\defref{c21-pick}\); subtree sizes from \(\chref{c24}\)). Use it to feel that
     ``unicyclic graph'' does \emph{not} mean offline queries.}
  \practice{Beautiful Subarrays}{CF 665E}
    {https://codeforces.com/problemset/problem/665/E}{CF 2100 L5}
    {Count subarrays with at most \(k\) even elements. The set-up is the closest call in
     the list -- offline, overlapping ranges over one array -- but the statistic is a
     count of evens, fully \emph{invertible}: a two-pointer window slides it in
     \(\bigO n\), a prefix count or BIT in \(\bigO{n\log n}\). Test~3 fails, so even
     \(\bigO{(n+q)\sqrt n}\) Mo is a tax on the array version of the lesson 27.2 taught
     with CF~383C (\(\defref{c21-signature}\)). Drill: say \emph{why} the window here
     never needs a block sort.}
  \practice{Timofey and a Black-White Tree}{CF 1790F}
    {https://codeforces.com/problemset/problem/1790/F}{CF 2100 L5}
    {Nodes start white and turn black one by one; at time \(i\), before recolouring
     node~\(i\), report how many black nodes are reachable from \(i\) through black
     nodes only. This is component \emph{size} under a growing activation prefix --
     connectivity, not a reorderable range: the answers must respect the time prefix
     (test~1 fails). The standard move is reverse-time DSU merging black nodes for the
     late times plus a thresholded direct DFS for the early ones (sqrt-decomposition
     over \emph{time}, not over ranges). Keep the DSU instinct from
     \(\secref{c21-tarjan}\); reject Mo's block order (\(\defref{c21-pick}\)).}
  \practice{We Need More Bosses}{CF 1000E}
    {https://codeforces.com/problemset/problem/1000/E}{CF 2100 L5}
    {Query \((s,t)\): the maximum number of bridges any walk from \(s\) to \(t\) can
     force. Compress the 2-edge-connected components; what remains is a \emph{bridge
     tree}, and the answer is just the number of tree-edges on the \(s\)-\(t\) path --
     a depth-prefix sum with an \(\lca\) correction. The distance decomposes over
     prefixes (test~3 fails) and each query answers online: reject Mo, keep bridges
     \(+\) lifting (\(\chref{c11}\)). The compression step is the chapter's one
     inversion: run Tarjan-family DFS ideas \emph{for} the structure, not against it.}
  \practice{Tree Master}{CF 1806E}
    {https://codeforces.com/problemset/problem/1806/E}{CF 2200 L5}
    {Tags \emph{brute force, data structures, dfs, trees}. Feed it to the four tests
     before coding: is the whole query set known up front, do the queried sets overlap,
     does the statistic refuse to telescope, and is one element cheap to add and drop?
     A \emph{data structures + dfs} tree problem with a \emph{brute force} tag is the
     tell for ``a fold that a smarter structure beats'' -- so first look for the
     invertible/structural route (Fenwick over the flatten, or DSU-on-tree for
     subtree-all, \(\chref{c28}\)); only if all four tests fire does Mo earn its
     \(\sqrt n\) (\(\defref{c21-signature}\), \(\secref{c21-boundary}\)).}
  \practice{Weights Division (hard version)}{CF 1399E2}
    {https://codeforces.com/problemset/problem/1399/E2}{CF 2200 L5}
    {Tags \emph{binary search, dfs and similar, greedy, sortings, trees}. Weighted
     tree rooted at \(1\), threshold \(S\): repeatedly halve an edge (floor), each
     halving costs \(1\), and every root-to-leaf sum must end up \(\le S\); minimise
     the operations. There is \emph{no query list} to sweep (test~2 empty) and no
     window statistic at all: the leverage of halving edge \(e\) is
     \(\mathit{leaves}(e)\cdot(w-\lfloor w/2\rfloor)\) -- a bottom-up leaf-count
     fold plus a greedy/priority queue (the tree-DP family, \(\chref{c10}\)).
     Reject Mo the way 27.5 rejects every fold-shaped objective
     (\(\defref{c21-pick}\)); the \(\sqrt n\) has nothing to amortise here.}
  \practice{Yasya and the Mysterious Tree}{CF 1980G}
    {https://codeforces.com/problemset/problem/1980/G}{CF 2300 L5}
    {Tags \emph{bitmasks, data structures, dfs, graphs, greedy, strings, trees}. Path
     questions where the state per edge is a small \emph{parity bitmask}. A bitmask fold
     along \(u\)-to-\(v\) is invertible (xor of root-to-node masks, \(\lca\) correction),
     so a single path query is answered in \(O(\log)\) from prefixes -- the opposite of
     the ``no inverse'' Mo regime. Gate test~3 fails; use LCA + bitmask prefix
     (\(\chref{c11}\)). Reach for Mo only if the statistic were a distinct-count, which
     it is not (\(\defref{c21-signature}\)).}
\end{mixedset}

\chaptersection{The Summit -- CF 1585E Frequency Queries}

\begin{notebox}[Why this is the chapter top rope]
Every problem before it put a \emph{histogram} on a subtree or a path and let Mo sweep.
The Summit asks for something finer: on the path \(v\to\) root, look at the
\emph{frequency of each value}, keep the values seen at least \(l\) times, and report the
\(k\)-th one ordered \emph{by frequency}. That is a statistic on a statistic -- it has no
inverse and no single running scalar, so naive Mo (which carries one number per window)
does not fit. The rescue is the deepest idea of the chapter wearing a new hat: a plain
DFS \emph{already} walks the root-to-\(v\) path as a stack, pushing a value on entry and
popping it on exit -- the exact \(O(1)\) \texttt{add}/\texttt{rm} of \(\secref{c21-moengine}\),
now riding the recursion instead of two pointers. Layer a frequency-of-frequency Fenwick
on top and the \(k\)-th-by-count is a select in \(O(\log n)\). It reuses \textbf{two}
tools of this chapter (the Mo frequency table and Tarjan's ``answer offline queries
attached to a node during one DFS'') plus \textbf{one} from Ch.~5 (the depth/parent DFS),
and is hand-replayed to both judge outputs \texttt{1 -1 1} and \texttt{1 1 2 1 -1}.
\end{notebox}

\solvedproblem{c21-summit}{Frequency Queries}{CF 1585E}
  {https://codeforces.com/problemset/problem/1585/E}{CF 2400 L6}
  {one DFS from the root turns the recursion stack into the Mo window (push a value on
   entry, pop on exit -- \(O(1)\) add/remove); maintain \(\mathit{cnt}[x]\) = occurrences
   of value \(x\) on the current root path and a frequency-of-frequency structure
   \(\mathit{bucket}[f]\) = number of distinct values seen exactly \(f\) times, kept as a
   Fenwick so ``\(k\)-th distinct value by increasing frequency among those with freq
   \(\ge l\)' ' is a prefix-sum select \(O(\log n)\); a per-frequency value stack returns
   the answer value; queries are attached to their node and answered at that DFS instant
   (Tarjan's offline discipline)}

\begin{problemstatement}[Frequency Queries (statement, constraints and both samples verified against the live CF page)]
Vertex \(1\) roots a tree on \(n\) nodes; node \(i\) carries an integer \(a_i\)
(\(1\le a_i\le n\)). There are \(q\) questions; question \((v,l,k)\) is answered by:
write the sequence of values on the shortest path from \(v\) up to the root (both
included); count how many times each value occurs; \textbf{drop} every value occurring
fewer than \(l\) times; \textbf{deduplicate} and \textbf{order the remaining distinct
values by their occurrence count in the original path sequence, increasing} (ties in any
order); the answer is the \(k\)-th value in that ordered list, or \(-1\) if fewer than
\(k\) values survive.

\textbf{Input.} Multiple test cases, total \(\sum n,\sum q\le 10^6\), \(1\le t\le10^6\).
Per case: \(n,q\) (\(1\le n,q\le10^6\)); the \(n\) values \(a_1..a_n\); the \(n-1\)
parents \(p_2..p_n\); then \(q\) lines \(v,l,k\) (\(1\le v,l,k\le n\)).

\textbf{Output.} For each question print its answer; ties may print any valid value.

\textbf{Samples (both official).} Case~1: \(n{=}3,q{=}3\), \(a{=}[1,1,1]\), parents
\([1,2]\) (chain \(1\!-\!2\!-\!3\)), questions \((3,1,1),(3,1,2),(3,2,1)\) \(\to\)
\texttt{1 -1 1}. Case~2: \(n{=}5,q{=}5\), \(a{=}[1,2,1,1,2]\), parents \([1,1,2,2]\),
questions \((3,1,1),(2,1,2),(4,1,1),(4,2,1),(4,2,2)\) \(\to\) \texttt{1 1 2 1 -1}.
\end{problemstatement}

\begin{attemptgate}[Attempt it yourself first -- \lv{6}, budget a real sitting]
The queried set is always a root path, and you need the whole frequency \emph{histogram}
of that path, not one scalar.  Why does that break plain Mo, and what does a DFS already
give you for free?
\peek{a root-to-\(v\) path is exactly the node-stack of a DFS when the search is at \(v\).
Enter a node: its value's count goes up by one. Leave: down by one. That is an \(O(1)\)
\texttt{add}/\texttt{rm} -- the Mo primitive -- but the ``window'' is the recursion depth,
so you need no block order at all.}
\nudge{Attach every question to its node \(v\) and answer it the instant the DFS is at
\(v\) (Tarjan's offline discipline, \(\secref{c21-tarjan}\)). Maintain
\(\mathit{cnt}[\mathit{val}]\) and \(\mathit{bucket}[f]=\#\{\text{values with }\mathit{cnt}=f\}\).
The ordered list you must select from is ``values with \(\mathit{cnt}\ge l\), by
\(\mathit{cnt}\) up'' -- a selection over the \(\mathit{bucket}\) array.}
\fullhint{Put \(\mathit{bucket}[\cdot]\) in a Fenwick so its prefix sum
\(S(f)=\sum_{g\le f}\mathit{bucket}[g]\) is ``\#distinct values with freq \(\le f\)'' and
advances by \(\pm1\) whenever a value moves between frequencies. For \((v,l,k)\):
\(\mathit{base}=S(l-1)\), \(\mathit{total}=S(n)\); if \(\mathit{total}-\mathit{base}<k\)
answer \(-1\); else find the smallest \(f\) with \(S(f)\ge \mathit{base}+k\) (Fenwick
lower-bound, \(O(\log n)\)) and output \emph{any} value whose current count is \(f\) --
maintain a per-frequency swap-remove stack \(\mathit{valAt}[f]\) to read one in \(O(1)\).
Moving value \(x\) from count \(c\) to \(c{+}1\): \(\mathit{bit}.\mathrm{add}(c,-1),
\mathit{bit}.\mathrm{add}(c{+}1,+1)\) and relocate \(x\) between \(\mathit{valAt}[c]\)
and \(\mathit{valAt}[c{+}1]\) (skip the \(c=0\) add). On leaving, mirror it with \(c\to
c{-}1\). One DFS answers all \(q\) in \(\bigO{(n+q)\log n}\); the structures fully unwind
to empty when the DFS returns, so the next test case starts clean.}
\end{attemptgate}

\solstep{Understanding}
Two reframings collapse the problem. First, ``the sequence on the path \(v\to\) root'' is
not an arbitrary range needing a doubled tour -- it is the \emph{ancestor stack} the DFS
holds while it sits at \(v\). So instead of Mo's two pointers creeping over a static
array, the window is the current recursion path, and it is entered/exited one node at a
time for free: entering \(x\) does \(\mathit{cnt}[a_x]{+}{+}\), leaving does
\(\mathit{cnt}[a_x]{-}{-}\), each an \(O(1)\) add/remove of exactly the \(\secref{c21-moengine}\)
flavour. Second, the answer is a \emph{selection on the frequency histogram}: among
distinct values with \(\mathit{cnt}\ge l\), ordered by \(\mathit{cnt}\) ascending, take
the \(k\)-th. That histogram is captured by \(\mathit{bucket}[f]\) -- ``how many distinct
values occur exactly \(f\) times'' -- and the selection is a prefix-sum lower-bound over
\(\mathit{bucket}\). A Fenwick maintains \(\mathit{bucket}\) under \(\pm1\) moves and
answers the select in \(O(\log n)\); the offline ``answer each question at its node''
discipline is Tarjan's (\(\secref{c21-tarjan}\)), not Mo's block order. The Summit is
therefore the chapter's thesis -- cheap single-element updates plus a non-invertible
histogram -- solved by the \emph{other} window mechanism this chapter introduced.

\solstep{Key Observations}
\begin{observation}
\textbf{A root path is a DFS stack.} While the DFS is at \(v\), the multiset of values
with \(\mathit{cnt}>0\) is exactly the values on the path \(v\to\) root. Entering/leaving
a node changes \(\mathit{cnt}\) of one value by one, so the window updates are \(O(1)\)
and come in matched enter/leave pairs -- everything unwinds to empty at the end of the
DFS, which is why no explicit reset is needed between test cases.
\end{observation}
\begin{observation}
\textbf{Ordering by frequency is selection on \(\mathit{bucket}\).} The deduplicated,
freq-increasing list is the values grouped by their count \(f=l,l+1,\dots\), group \(f\)
holding \(\mathit{bucket}[f]\) values. The \(k\)-th survivor is any value in the group
\(f^\ast\) where \(f^\ast\) is the smallest \(f\ge l\) with
\(\sum_{g=l}^{f}\mathit{bucket}[g]\ge k\). With \(S(f)=\sum_{g\le f}\mathit{bucket}[g]\)
that is ``smallest \(f\) with \(S(f)\ge S(l-1)+k\)'': a Fenwick lower-bound in \(O(\log n)\),
and ties let us emit any member of group \(f^\ast\) (read from \(\mathit{valAt}[f^\ast]\)).
\end{observation}

\solstep{Think It With Me}
\begin{thinkbox}
\deadend{Run Mo on paths: treat the root path as a doubled-tour window}
{Path Mo (\(\secref{c21-path}\)) answers a statistic that is a single running scalar
(distinct count, \(\sum\mathrm{freq}^2\)) -- you read one number when the window is right.
Here the answer needs the \emph{entire} frequency histogram re-sorted and a \(k\)-th
element selected, which Mo's scalar window cannot express; you would re-scan all
frequencies per query (\(\bigO{n}\) each, \(\tle\)). And a root path is already
contiguous in the ancestor stack -- the doubled-tour toggle is needless overhead.
\textbf{TLE / wrong model}.}
\deadend{Keep \(\mathit{bucket}[f]\) in a plain array and scan \(f=l,l+1,\dots\) per query}
{Each query then costs \(O(\text{number of distinct frequencies})=O(n)\) in the worst
case, and \(q\le 10^6\), so \(O(nq)\). The cumulative-count selection over
\(\mathit{bucket}\) is a \emph{prefix-sum lower-bound} -- exactly what a Fenwick does in
\(O(\log n)\). You must promote the histogram array to a searchable structure.
\textbf{TLE}.}
\deadend{Recompute the path value multiset for every query from \(v\) up through parents}
{Walking \(v\to\) root per query is \(\bigO{\depth}\) and rebuilding a frequency map from
scratch kills \(q\) -- no work is shared between sibling queries. The whole point of the
DFS view is that moving from a node to its child changes the histogram by \emph{one}
value; re-walking per query throws that amortisation away. \textbf{TLE}.}
\aha{Do one DFS from the root and let the recursion stack be the window: \(\mathit{enter}(x)\)
bumps \(\mathit{cnt}[a_x]\) (moving value \(a_x\) from frequency bucket to bucket in a
Fenwick and between two \(\mathit{valAt}\) stacks) and \(\mathit{leave}(x)\) undoes it.
Answer every query attached to \(v\) at the DFS instant: \(\mathit{base}=S(l-1)\),
\(\mathit{total}=S(n)\), if \(\mathit{total}-\mathit{base}<k\) print \(-1\), else
Fenwick-lower-bound for the smallest \(f\) with \(S(f)\ge\mathit{base}+k\) and print any
member of \(\mathit{valAt}[f]\). All in \(\bigO{(n+q)\log n}\).}
\end{thinkbox}

\solstep{The Approach Ladder}
\begin{center}\renewcommand{\arraystretch}{1.25}\setlength{\tabcolsep}{3pt}
\begin{tabular}{@{}c p{0.23\linewidth} p{0.29\linewidth} l c@{}}
\toprule
\textbf{\#} & \textbf{Approach} & \textbf{Idea} & \textbf{Time} & \textbf{Verdict}\\
\midrule
1 & \lv{3}\ walk \(v\to\) root per query & rebuild freq map from scratch & \(\bigO{q\cdot n}\) & \tle \\
2 & \lv{4}\ path Mo, doubled tour & scalar window, wrong shape & \(\bigO{(n+q)\sqrt n}\) & \textcolor{warningred}{\ding{55}\ \textbf{WA}} \\
3 & \lv{5}\ DFS stack window + array histogram & right window, linear select & \(\bigO{(n+q)\cdot n}\) & \tle \\
4 & \lv{6}\ \textbf{DFS stack + freq-of-freq Fenwick} & push/pop + lower-bound select & \(\bigO{(n+q)\log n}\) & \ac \\
\bottomrule
\end{tabular}\end{center}
Rung~4 is the Summit: rungs~1--3 each keep the right \emph{goal} but the wrong
\emph{object} (a re-walked path, a scalar Mo window, an unsearchable histogram); realising
that a root path \emph{is} a DFS stack and that ``\(k\)-th by frequency'' \emph{is} a
Fenwick lower-bound over the frequency-of-frequency array is the whole win.

\solstep{Solution}
Build child lists from the parent array. DFS from node \(1\): on \(\mathit{enter}(u)\)
move value \(a_u\) from its current frequency to the next (\(\mathit{cnt}\) up), applying
the matched \(\mathit{bucket}\) Fenwick \(-1/+1\) and relocating \(a_u\) between
\(\mathit{valAt}\) stacks; then answer every query attached to \(u\) by
\(\mathit{base}=S(l-1),\ \mathit{total}=S(n)\), emitting \(-1\) if
\(\mathit{total}-\mathit{base}<k\) else the value \(=\mathit{valAt}[f^\ast].\!\)back where
\(f^\ast=\)Fenwick-lower-bound\((\mathit{base}+k)\). Recurse into children, then
\(\mathit{leave}(u)\) mirrors the move. Because every enter is matched by a leave, all
structures return to empty when the DFS ends, so per test case you only clear the
adjacency and query lists. Read answers in input order.

\solstep{Proof of Correctness}
\begin{proofbox}
\textbf{(Window $=$ path.)} At the DFS instant at \(v\), the entered-but-not-yet-left nodes
are exactly the root-to-\(v\) stack, so \(\{x:\mathit{cnt}[x]>0\}\) equals the multiset of
values on the path \(v\to\) root; enter/leave maintain this by induction on DFS events.
\textbf{(Histogram correctness.)} \(\mathit{bucket}[f]=|\{x:\mathit{cnt}[x]=f\}|\) is
updated \(-1\) on the old count and \(+1\) on the new whenever a single value moves, so it
always reflects the current path's frequency-of-frequency, and the Fenwick stores it, giving
\(S(f)=\#\{x:\mathit{cnt}[x]\le f\}\). \textbf{(Selection.)} The required ordered list is the
values with \(\mathit{cnt}\ge l\), ascending by count: group \(l\) first, then \(l{+}1\), etc.,
each group contributing \(\mathit{bucket}\) members in arbitrary (tie) order. The \(k\)-th
element lies in the smallest group \(f^\ast\ge l\) with \(\sum_{g=l}^{f^\ast}\mathit{bucket}[g]\ge k\),
i.e.\ \(S(f^\ast)\ge S(l-1)+k\) and \(S(f^\ast{-}1)<S(l-1)+k\) -- the Fenwick lower-bound of
\(\mathit{base}+k\). Existence fails exactly when \(S(n)-\mathit{base}<k\), giving \(-1\);
otherwise any member of group \(f^\ast\) (top of \(\mathit{valAt}[f^\ast]\), nonempty because
\(\mathit{bucket}[f^\ast]\ge1\)) is a valid answer. \textbf{(Unwind.)} Each node's enter is
mirrored by its leave, so after \(\mathit{dfs}(1)\) all counts, buckets and stacks are empty:
no inter-test contamination. Hence every printed answer matches the specification. \(\Box\)
\end{proofbox}

\solstep{Implementation}
\begin{lstlisting}[caption={CF 1585E Summit: a DFS from the root uses the recursion stack as the Mo window (enter bumps a value's frequency, leave undoes it); a frequency-of-frequency Fenwick selects the k-th distinct value by increasing count, and per-frequency value stacks hand back an answer value. O((n+q) log n). ASCII only, identifiers without underscores.}, label={lst:c21-summit}]
#include <bits/stdc++.h>
using namespace std;
const int MAXN=1000005;
int T,n,q,highest,a[MAXN];
vector<int> kid[MAXN];
struct QQ{ int l,k,id; };
vector<QQ> atNode[MAXN];
int ansArr[MAXN];
int cntVal[MAXN];            // occurrences of value on current root path
int bitFen[MAXN];            // Fenwick over bucket[f] = #values with freq f
vector<int> valAt[MAXN];     // the values currently having frequency f
int posIn[MAXN];             // index of a value inside valAt[its frequency]
void bitAdd(int i,int d){ for(;i<=n;i+=i&-i) bitFen[i]+=d; }
int bitSum(int i){ int s=0; for(;i>0;i-=i&-i) s+=bitFen[i]; return s; }
int bitLower(int target){ int pos=0;
  for(int step=highest; step; step>>=1){ int np=pos+step;
    if(np<=n && bitFen[np]<target){ pos=np; target-=bitFen[np]; } }
  return pos+1; }
void attachVal(int v,int f){ valAt[f].push_back(v); posIn[v]=valAt[f].size()-1; }
void detachVal(int v,int f){ int p=posIn[v]; int last=valAt[f].back();
  valAt[f][p]=last; posIn[last]=p; valAt[f].pop_back(); }
void moveVal(int v,int from,int to){         // value v changes count from->to
  detachVal(v,from); bitAdd(from,-1);
  attachVal(v,to);   bitAdd(to,1); }
void enter(int u){ int v=a[u], c=cntVal[v];
  if(c>0) moveVal(v,c,c+1); else { attachVal(v,1); bitAdd(1,1); }
  cntVal[v]=c+1; }
void leave(int u){ int v=a[u], c=cntVal[v];
  if(c>1) moveVal(v,c,c-1); else { detachVal(v,1); bitAdd(1,-1); }
  cntVal[v]=c-1; }
void dfs(int u){
  enter(u);
  for(const QQ &Q : atNode[u]){
    int base=bitSum(Q.l-1), total=bitSum(n);
    if(total-base < Q.k) ansArr[Q.id]=-1;
    else { int f=bitLower(base+Q.k); ansArr[Q.id]=valAt[f].back(); }
  }
  for(int w : kid[u]) dfs(w);
  leave(u);
}
int main(){
  scanf("%d",&T);
  while(T--){
    scanf("%d%d",&n,&q);
    for(int i=1;i<=n;i++) scanf("%d",&a[i]);
    for(int i=2;i<=n;i++){ int p; scanf("%d",&p); kid[p].push_back(i); }
    for(int i=1;i<=q;i++){ int v,l,k; scanf("%d%d%d",&v,&l,&k); atNode[v].push_back({l,k,i}); }
    highest=1; while((highest<<1)<=n) highest<<=1;
    dfs(1);
    for(int i=1;i<=q;i++) printf("%d ",ansArr[i]);
    printf("\n");
    for(int i=1;i<=n;i++){ kid[i].clear(); atNode[i].clear(); } // window auto-unwinds
  }
}
\end{lstlisting}

\solstep{Trace on the Official Samples}
\begin{tracebox}[Sample 1: chain \(1\!-\!2\!-\!3\), all values \(1\) -- the histogram is one spike]
At node~3 the root path is \(1,2,3\), values \([1,1,1]\), so \(\mathit{cnt}[1]=3\) and
\(\mathit{bucket}[3]=1\) (S(1)=S(2)=0, S(3)=1). \((3,1,1)\): base=S(0)=0, total=S(3)=1,
\(1{-}0\ge1\), target=1, bitLower(1)=3, valAt[3]=\{1\} \(\Rightarrow\) \(\boxed{1}\).
\((3,1,2)\): total\(-\)base\(=1<2\) \(\Rightarrow\) \(\boxed{-1}\). \((3,2,1)\): base=S(1)=0,
total=1, target=1 \(\Rightarrow f^\ast=3\) \(\Rightarrow\) \(\boxed{1}\). Output
\texttt{1 -1 1} -- match. \(\Box\)
\end{tracebox}

\begin{tracebox}[Sample 2: the trap that ties decide, and the LCA-instant histogram]
Tree \(1\!\to\!\{2,3\}\), \(2\!\to\!\{4,5\}\); \(a=[1,2,1,1,2]\).
\((3,1,1)\) at node 3, path values \([1,1]\): \(\mathit{cnt}[1]=2\), bucket\([2]=1\),
S(0)=0, total=S(5)=1, target=1, \(f^\ast=2\), val=1 \(\Rightarrow\boxed{1}\).
\((2,1,2)\) at node 2, path \([1,2]\): \(\mathit{cnt}[1]=\mathit{cnt}[2]=1\),
bucket\([1]=2\), base=S(0)=0, total=2, target=2, \(f^\ast=1\) (S(1)=2\(\ge\)2),
valAt[1]=\(\{1,2\}\) \(\Rightarrow\) \(\boxed{1\text{ (or }2\text{)}}\), a legal tie pick.
\((4,1,1)\) at node 4, path \([1,2,1]\): \(\mathit{cnt}[1]=2,\mathit{cnt}[2]=1\),
bucket\([1]=1,\)bucket\([2]=1\), target=1 \(\Rightarrow f^\ast=1\), val=2 \(\Rightarrow\boxed{2}\).
\((4,2,1)\): base=S(1)=1, total=2, target=2, \(f^\ast=2\), val=1 \(\Rightarrow\boxed{1}\).
\((4,2,2)\): total\(-\)base\(=2-1=1<2\) \(\Rightarrow\boxed{-1}\). Output
\texttt{1 1 2 1 -1} (the tie on query~2 printed as 1) -- match. \(\Box\)
\end{tracebox}

\solstep{Complexity}
\begin{complexitybox}
Each node is entered and left once; every move does two Fenwick updates and one
\(\mathit{valAt}\) relocate, \(O(\log n)\) amortised. Each query is answered with two
Fenwick sums and one lower-bound, \(O(\log n)\). Total \(\bigO{(n+q)\log n}\) per test,
\(\bigO{(\sum n+\sum q)\log n}\) over all cases -- inside the 4\,s limit at
\(\sum n,\sum q\le10^6\). \textbf{Space} \(\bigO{n+q}\): the child/query lists, the
Fenwick, \(\mathit{cnt}\), and the per-frequency \(\mathit{valAt}\) stacks whose total
size never exceeds the number of distinct values on the current path (\(\le n\)). No node
is ever re-walked and no histogram rebuilt -- the DFS-stack window and the
frequency-of-frequency Fenwick are the entire machinery.
\end{complexitybox}

\begin{hardeningbox}
\begin{itemize}[nosep,leftmargin=1.4em]
  \item \textbf{Do not run Mo here.} The answer is not a single running scalar; it needs
        the whole histogram re-sorted. Path Mo's window (\(\secref{c21-path}\)) gives the
        right \emph{set} but the wrong \emph{read}, and re-scanning frequencies per query
        is \(\tle\) (dead end 1 and 2).
  \item \textbf{The lower-bound lands in \(f^\ast\ge l\) for free.} Because
        \(\mathit{target}=\mathit{base}+k > S(l-1)\), the smallest \(f\) with
        \(S(f)\ge\mathit{target}\) is automatically \(\ge l\) -- you never special-case
        the ``drop counts below \(l\)'' step; it is absorbed by \(\mathit{base}=S(l-1)\).
  \item \textbf{The window auto-unwinds -- rely on it.} Matched enter/leave return every
        \(\mathit{cnt}\), bucket and \(\mathit{valAt}\) to empty at the end of the DFS, so
        per test you clear only adjacency and query lists; do \emph{not} memset the whole
        \(\bigO{n}\) arrays (the sum of \(n\) is \(10^6\), but a full-array clear per tiny
        test across \(t\le10^6\) cases would re-introduce an \(O(t\cdot n)\) blow-up).
  \item \textbf{Recursion depth at \(n=10^6\).} A chain overflows the call stack; convert
        \(\mathit{dfs}\) to an explicit stack (push enter, then children, then leave) --
        the same caution as Tarjan (\(\secref{c21-tarjan}\)).
\end{itemize}
\end{hardeningbox}

\begin{recallbox}[Chapter 27 in ten seconds]
\emph{When a huge offline batch of subtree/path queries has a statistic with no inverse,
block-sort the ranges and sweep one window with an \(O(1)\) add/remove --
\(\bigO{(n+q)\sqrt n}\) (\(\thmref{c21-complexity}\), \(\coderef{c21-mo}\)); a subtree is
the preorder range \([\tin(v),\tout(v)]\) (\(\thmref{c21-subtree-range}\)) and a path is a
doubled-tour toggle window split at the \(\lca\) plus the read-time LCA
(\(\thmref{c21-pathsplit}\), \(\coderef{c21-euler-toggle}\)); CF~617E is the array engine on
prefix xors (\(\probref{c21-xor}\)); CF~463E carries the root-path window on the DFS
stack with a per-prime push/undo table (\(\probref{c21-caisa}\)), while the
parity-normalised invertible case (Ch.~18's CF~383C, \(\chref{c14}\)) skips Mo
entirely; Tarjan's one DFS \(+\) DSU answers \(q\) LCAs in
\(\bigO{(n+q)\alpha}\) (\(\thmref{c21-tarjan}\)); and the four tests decide Mo versus
Fenwick / HLD / DSU-on-tree / virtual tree / centroid (\(\defref{c21-signature}\),
\(\defref{c21-pick}\)). The Summit CF~1585E (\(\probref{c21-summit}\)) turns the DFS stack
itself into the window and a frequency-of-frequency Fenwick into the selector.}
\selftest{
  \item State the four Mo tests and why CF~463E needs none of them. \hfill\lv{4}
  \item Give the doubled-tour window for path \((u,v)\) when \(u\) is not the ancestor, and the extra step. \hfill\lv{5}
  \item Why does Tarjan beat lifting only for \emph{offline} batches, and at what cost? \hfill\lv{4}
  \item Which structure turns ``\(k\)-th value by increasing frequency'' into one \(O(\log n)\) query? \hfill\lv{6}
}
\end{recallbox}

\chaptersection{Tool Kit Inventory}

The machine assembled through Chapter~27. Earlier chapters' rows are carried forward;
this chapter's are appended. Pointers only.

\begin{center}
\small
\renewcommand{\arraystretch}{1.25}
\setlength{\tabcolsep}{3pt}
\begin{tabular}{@{}p{0.24\linewidth} l c p{0.22\linewidth} >{\raggedright\arraybackslash\scriptsize}p{0.17\linewidth}@{}}
\toprule
\textbf{Tool} & \textbf{Since} & \textbf{Lv} & \textbf{Trigger phrase} & \textbf{Where} \\
\midrule
\multicolumn{5}{l}{\itshape Carried forward} \\
LCA + \(\dist\), binary lifting & 16 & \lv{4} & ``the split point of a path'' & \chref{c11}, \thmref{c11-distf} \\
Euler \(\tin/\tout\), flatten & 5 & \lv{3} & ``a subtree is one array range'' & \chref{c24} \\
Fenwick / difference array & 5/10 & \lv{3} & ``invertible range add, point query'' & \chref{c24}, \chref{c10} \\
DSU (find/union, path compress) & 17 & \lv{3} & ``merge finished regions'' & \chref{c28} \\
DSU on tree / small-to-large & 17 & \lv{4} & ``subtree statistic for (almost) every node'' & \chref{c28} \\
\midrule
\multicolumn{5}{l}{\itshape Chapter 27 -- new this chapter} \\
Mo block order + two pointers & 27.1 & \lv{4} & ``offline ranges, non-invertible \(O(1)\) update'' & \defref{c21-mo}, \coderef{c21-mo} \\
Prefix-xor pair counting & 27.1 & \lv{4} & ``count subarrays with \(\mathit{xor}=k\)'' & \probref{c21-xor}, \coderef{c21-xor} \\
Subtree flatten \(+\) Mo & 27.2 & \lv{4} & ``aggregate over subtree\((v)\), offline'' & \thmref{c21-subtree-range}, \coderef{c21-flatten} \\
Root-path stack + last-seen table & 27.2 & \lv{4} & ``deepest ancestor sharing a key'' & \probref{c21-caisa}, \coderef{c21-caisa} \\
Doubled-tour toggle + LCA split & 27.3 & \lv{5} & ``statistic on path \(u\)-\(v\), offline'' & \thmref{c21-pathsplit}, \coderef{c21-euler-toggle} \\
Tarjan offline LCA & 27.4 & \lv{4} & ``answer \(q\) LCAs, no lifting table'' & \thmref{c21-tarjan}, \coderef{c21-tarjan} \\
Recognition gate / boundary pick & 27.4/27.5 & \lv{5} & ``is this even a Mo job?'' & \defref{c21-signature}, \defref{c21-pick} \\
Mo kit (decision skeleton) & 27 & \lv{5} & ``the seven questions in one place'' & \coderef{c21-kit} \\
Summit: DFS-stack window + freq-of-freq BIT & Summit & \lv{6} & ``root path, order values by count'' & \coderef{c21-summit} \\
\bottomrule
\end{tabular}
\end{center}

\chaptersection{Related Topics}

\begin{itemize}[nosep, leftmargin=1.4em]
  \item \chref{c30} (Ch.~26, Virtual Trees) --- backward, \textbf{the hand-off this chapter
        answers}: Ch.~26's boundary (\(\secref{c30-boundary}\)) forwarded ``many overlapping
        offline node-set queries with a non-invertible statistic'' here; \(\secref{c21-moengine}\)
        opens by discharging exactly that debt. The virtual tree re-folds \emph{per marked
        set}; Mo \emph{shares work across} the batch.
  \item \chref{c11} (Ch.~16, LCA) --- tool: \(\lca\) splits the path window (\(\secref{c21-path}\))
        and is the quantity Tarjan (\(\secref{c21-tarjan}\)) computes offline; \(\thmref{c11-distf}\)
        is used, never re-derived.
  \item \chref{c24} (Ch.~5, Euler Tour / \(\tin/\tout\)) --- tool: the subtree flatten
        (\(\defref{c21-flatten}\)) is Ch.~5's array; the doubled tour of \(\secref{c21-path}\)
        extends it, and Ch.~18's difference-Fenwick solve of CF~383C
        (\(\chref{c14}\)) is Ch.~5's invertible primitive.
  \item \chref{c28} (Ch.~17, DSU on tree) --- contrast and boundary: DSU-on-tree and Mo both
        answer subtree distinct-colour folds (CF~375D, whose URL lives there); \(\secref{c21-boundary}\)
        prices \(\bigO{n\log n}\) DSU against \(\bigO{(n+q)\sqrt n}\) Mo and picks per shape.
  \item \chref{c16} (Ch.~22, HLD) / \chref{c23} (Ch.~25, Link-Cut) --- boundary: every
        \emph{online} path or update-interleaved query leaves Mo (which is offline-only)
        for these (\(\defref{c21-signature}\)).
  \item \chref{c18} (Ch.~24, Centroid Decomposition) --- boundary: whole-tree pair/distance
        distributions with no query list go to centroid, not a window sweep
        (\(\secref{c21-boundary}\)); the Summit's frequency-of-frequency also echoes centroid's
        ``fold over a global distribution''.
  \item \chref{c22} (Ch.~28, Tree Isomorphism and Hashing) --- \textbf{next chapter}: where Mo
        compares \emph{counts}, hashing compares \emph{structures}; the note in
        \(\secref{c21-boundary}\) hands the ``which subtree equals which'' question forward.
\end{itemize}

\begin{lastword}[Last Word --- Chapter 27]
\emph{``When a statistic has no inverse, do not try to be clever about prefixes -- be
clever about \textbf{order}. Block-sort the offline queries so consecutive windows differ
by one node, keep an \(O(1)\) add/remove, and a \(\bigO{nq}\) nightmare becomes
\(\bigO{(n+q)\sqrt n}\); the tree enters through two doors, the preorder flatten (a subtree
is one range) and the doubled Euler tour (a path is a toggle window split at the
\(\lca\)). And when the batch only wants \(\lca\), Tarjan's DFS \(+\) DSU answers it in
near-linear time with no lifting table at all.''}

\medskip
Chapter~26 handed this chapter a problem it could not compress, and Chapter~27 collected
it. Section 27.1 built the engine and proved \(\bigO{(n+q)\sqrt n}\)
(\(\thmref{c21-complexity}\)), cashing it in on CF~617E's prefix-xor pairs
(\(\probref{c21-xor}\)); 27.2 bolted on the tree via the subtree flatten
(\(\thmref{c21-subtree-range}\)), handed the scary depth-parity update to Ch.~18's
Fenwick as the invertible case it is (\(\chref{c14}\)), and opened a third door with
CF~463E -- the recursion stack carrying a per-prime push/undo table through
fifty-one rebuilds (\(\probref{c21-caisa}\)); 27.3 invented the \term{toggle} so a path -- which no
flatten can make contiguous -- becomes a window, and disarmed the LCA-parity trap by hand
on three walks of \(R\) (\(\thmref{c21-pathsplit}\)); 27.4 gave the offline \(\lca\) with
Tarjan's WHITE/GREY/BLACK DSU run traced to \(1,2,2,3,1\) (\(\thmref{c21-tarjan}\)); and
27.5 drew the boundary so you reach for the sweep only after Fenwick, HLD, DSU-on-tree,
virtual tree and centroid have each been offered the job and declined it
(\(\defref{c21-pick}\)). The Summit CF~1585E (\(\probref{c21-summit}\)) fused it all: the
DFS stack \emph{is} the Mo window, and a frequency-of-frequency Fenwick reads the answer.

\medskip
Carry three reflexes forward:
\begin{enumerate}[nosep, leftmargin=1.6em]
  \item \emph{Look for the inverse first.} If \(\mathrm{pref}(r)-\mathrm{pref}(l-1)\)
        exists -- even after a \((-1)^{\depth}\) normalisation -- Mo is a \(\sqrt n\) tax
        you chose to pay (\(\chref{c14}\)).
  \item \emph{Choose the array before the algorithm.} Subtree \(\to\) preorder flatten;
        path \(\to\) doubled tour with a toggle; neither is ``the tree'' -- the sweep runs
        on a line (\(\thmref{c21-subtree-range}\), \(\thmref{c21-pathsplit}\)).
  \item \emph{Respect modality.} Mo is offline and update-free; one interleaved update and
        the block order is illegal -- that is HLD's or the LCT's contract, not yours
        (\(\defref{c21-signature}\)).
\end{enumerate}

\medskip
And three habits: get \(B\approx n/\sqrt q\) right and verify the \texttt{add}/\texttt{rm}
are true inverses before trusting the sweep; add the path-Mo LCA at \emph{read} time only
in case (ii); and when a whole batch is LCA, reach for Tarjan (\(\bigO{(n+q)\alpha}\), no
table) before building an \(O(n\log n)\) lifting structure. You now own the one engine that
answers tree queries by \emph{reordering the questions} rather than by folding the tree --
the offline, amortised companion to every prefix, path, and per-set structure before it.

\medskip
\textbf{Next:} \chref{c22} (Tree Isomorphism and Hashing) changes the question entirely.
Mo and Tarjan compare \emph{counts} on sets of nodes; the next chapter asks when two trees
are \emph{the same tree} -- canonical forms (AHU), subtree hashing for structural equality,
centre-based unrooted comparison, and counting distinct subtrees -- the equality test this
chapter never needed, and the tool behind cutting-point and symmetry problems like CF~1800G
and CF~1252F.
\end{lastword}
